\documentclass[11pt]{article}

\usepackage[a4paper,margin=1.15in]{geometry}
\usepackage{amsmath,amssymb,amsthm,mathtools}
\usepackage{microtype}
\usepackage[hidelinks]{hyperref}
\usepackage[nameinlink,capitalise]{cleveref}

\newtheorem{theorem}{Theorem}[section]
\newtheorem{corollary}[theorem]{Corollary}
\newtheorem{proposition}[theorem]{Proposition}
\newtheorem{lemma}[theorem]{Lemma}
\theoremstyle{remark}
\newtheorem{remark}[theorem]{Remark}

\newcommand{\T}{\mathbb T}
\newcommand{\R}{\mathbb R}
\newcommand{\Z}{\mathbb Z}
\newcommand{\e}{\mathrm e}
\newcommand{\supp}{\operatorname{supp}}
\newcommand{\cl}{\operatorname{cl}}
\newcommand{\Rel}{\mathcal R}

% Supporting manuscript only; rejected as the flagship contribution on 2026-07-22.
\title{Compact Hulls, Parity Obstructions, and Value Recurrence\\
for Suzuki's Riemann Screw Function}
\author{Draft for priority and specialist review}
\date{22 July 2026}

\begin{document}
\maketitle

\begin{abstract}
Let $\Psi$ be Suzuki's function associated with the Riemann xi-function.  Suzuki proved that the
Riemann hypothesis is equivalent to boundedness of $\Psi$ and observed that, even under the
hypothesis, its limit inferior and limit superior are difficult to determine.  The endpoints he
displayed were obtained under the additional assumptions that the zeros are simple and their
positive ordinates are linearly independent over $\mathbb Q$.

We prove a compact-hull law for every absolutely summable positive cosine spectrum.  Its cofinal
value set is a closed interval, every approximate value recurs with bounded gaps, every interior
level is attained with bounded gaps, and its limiting value distribution is the pushforward of Haar
measure on the frequency-relation hull.  No independence hypothesis is required.  The absolute
upper endpoint is attained if and only if every finite integer relation among the frequencies has
even total coefficient.  This parity criterion is necessary and sufficient and is strictly weaker
than rational independence.

Applied to Suzuki's exact zero expansion, the theorem proves under the Riemann hypothesis, without
simplicity or linear independence, that $\liminf_{t\to\infty}\Psi(t)=0$, that every tail has the
same interval of limit values, and that a compactly supported limiting value distribution exists
with support equal to that interval.  It also yields the equivalence between the Riemann hypothesis
and uniform Bohr almost periodicity of $\Psi$.  The theorem does not prove the Riemann hypothesis;
it converts the remaining numerical $\limsup$ question into an explicit parity problem for integer
relations among zero ordinates.
\end{abstract}

\section{Introduction}

Write $\xi$ for the Riemann xi-function and enumerate, with multiplicity, the zeros $\gamma$ of
$z\mapsto\xi(1/2-iz)$.  Suzuki introduced the continuous even function
\begin{equation}\label{eq:suzuki-series}
  \Psi(t)=\sum_{\gamma}\frac{1-\cos(\gamma t)}{\gamma^2},
\end{equation}
where the identity is unconditional and the series is interpreted with the symmetries of the xi
divisor.  This function is the negative of Suzuki's Riemann screw function.  Among several exact
criteria, Suzuki proved
\begin{equation}\label{eq:suzuki-boundedness}
  \mathrm{RH}\quad\Longleftrightarrow\quad \Psi(t)=O(1)
  \qquad (t\geq 0)
\end{equation}
and, under RH, $\Psi(t)>0$ for $t\ne0$ \cite{Suzuki2023}.

Under RH all frequencies in \eqref{eq:suzuki-series} are real, so the series is an absolutely and
uniformly convergent positive cosine series.  Suzuki noted that its limit inferior and limit
superior are nevertheless not easy to determine.  Assuming in addition simplicity and rational
independence of the positive zero ordinates, he obtained
\begin{equation}\label{eq:suzuki-LI-endpoints}
 \liminf_{t\to\infty}\Psi(t)=0,
 \qquad
 \limsup_{t\to\infty}\Psi(t)=2\sum_\gamma\frac1{\gamma^2}.
\end{equation}
Later work of Matsumoto and Suzuki developed value-distribution and $M$-function results for
related zero sums, again using a linear-independence condition to construct a density
\cite{MatsumotoSuzuki2025}.  Suzuki's more recent finite-interval program instead studies the Weil
quadratic form and a conjectural spectral limit of self-adjoint differential realizations
\cite{Suzuki2026}.

The purpose of this paper is to remove independence from the recurrence and measure-existence
parts of the first problem.  The correct replacement for the full torus is the closed subgroup cut
out by all integer relations among the frequencies.  Three features then become visible.

First, the relation hull is connected because it is the closure of a one-parameter subgroup.
Consequently a real continuous observable on it has interval image: there are no gaps in the
cofinal value set.  Second, compact-group minimality upgrades subsequential recurrence to bounded-gap
recurrence.  Third, the inhomogeneous phase $(-1,-1,\ldots)$ belongs to the relation hull exactly
when every integer relation has even coefficient sum.  This is the precise obstruction to aligning
all cosine modes at their individual maxima.

Our universal theorem is stated in \cref{sec:universal}.  The full cofinal value and Haar laws are
proved in \cref{sec:recurrence}.  The parity theorem is \cref{thm:parity}.  In
\cref{sec:screw-hull} we identify the compact group as the intrinsic Hilbert hull of the associated
screw line.  The application to $\Psi$ is given in \cref{sec:xi}.

The harmonic-analysis ingredients are classical: compact abelian duality, Haar measure, and
Kronecker--Bohr recurrence.  The new claim requiring priority review is their complete
infinite-spectrum synthesis for Suzuki's exact screw series, including the no-independence
$\liminf$, support, syndetic-level, and parity conclusions.  We make no claim that this proves RH or
that it decides the integer-relation parity of the actual xi zeros.

\section{The relation hull of a summable screw spectrum}\label{sec:universal}

Let $I$ be a nonempty finite or countable set.  Fix distinct frequencies $\omega_i>0$ and weights
$a_i>0$ such that
\begin{equation}\label{eq:summable}
 A:=\sum_{i\in I}a_i<\infty.
\end{equation}
Define
\begin{equation}\label{eq:P-def}
 P(t):=\sum_{i\in I}a_i\bigl(1-\cos(\omega_i t)\bigr).
\end{equation}
Let $\T=\{z\in\mathbb C:|z|=1\}$, and consider the continuous homomorphism
\begin{equation}\label{eq:theta}
 \Theta:\R\longrightarrow\T^I,
 \qquad
 \Theta(t)_i=\e^{i\omega_i t}.
\end{equation}
Its relation hull is
\begin{equation}\label{eq:K-def}
 K:=\cl\Theta(\R)\subset\T^I.
\end{equation}
The finitely supported integer-relation lattice is
\begin{equation}\label{eq:relation-lattice}
 \Rel_\omega:=
 \left\{k\in\Z^{(I)}:\sum_{i\in I}k_i\omega_i=0\right\}.
\end{equation}

The positive observable on the hull is
\begin{equation}\label{eq:F-def}
 F(z):=\sum_{i\in I}a_i(1-\Re z_i).
\end{equation}
The Weierstrass test, with majorant $2a_i$, shows that \eqref{eq:P-def} and \eqref{eq:F-def}
converge uniformly.  In particular $F$ is continuous and
\begin{equation}\label{eq:P-factor}
 P=F\circ\Theta.
\end{equation}

\begin{theorem}[Relation-hull law]\label{thm:relation-hull}
The set $K$ is a compact connected abelian group, and
\begin{equation}\label{eq:annihilator-description}
 K=\left\{z\in\T^I:
       \prod_{i\in I}z_i^{k_i}=1\text{ for every }k\in\Rel_\omega\right\}.
\end{equation}
For every $T_0\in\R$, the positive tail $\Theta([T_0,\infty))$ is dense in $K$.
\end{theorem}

\begin{proof}
The product torus is compact and abelian.  Since $\Theta$ is a continuous homomorphism, $K$ is a
compact subgroup.  The image $\Theta(\R)$ is connected, and closure preserves connectedness, so
$K$ is connected.

Every continuous character of $\T^I$ is a finite-coordinate monomial
\[
 \chi_k(z)=\prod_{i\in I}z_i^{k_i},\qquad k\in\Z^{(I)}.
\]
On the one-parameter subgroup,
\[
 \chi_k(\Theta(t))=
 \exp\left(it\sum_{i\in I}k_i\omega_i\right).
\]
Thus $\chi_k$ is trivial on $K$ if and only if $k\in\Rel_\omega$.  A closed subgroup of a compact
abelian group is the intersection of the kernels of the characters trivial on it.  This proves
\eqref{eq:annihilator-description}.

It remains to check the direction of time, which is important for cofinal statements.  Put
\[
 S:=\cl\Theta([0,\infty))\subset K.
\]
This is a compact subsemigroup.  We show that it is a group.  Let $x\in S$.  Compactness gives a
subsequence $x^{n_j}$ converging to a point $y$.  After passing to a further subsequence,
\[
 x^{n_{j+1}-n_j}\longrightarrow 1.
\]
If the positive differences are eventually one, then $x=1$.  Otherwise take a subsequence on which
they are at least two; then
\[
 x^{n_{j+1}-n_j-1}\longrightarrow x^{-1}.
\]
The powers on the left lie in $S$, so $x^{-1}\in S$.  Hence $S$ is a group.  It contains the
nonnegative orbit and therefore its inverses, so it contains $\Theta(\R)$.  Thus $S=K$.
Multiplication by $\Theta(T_0)$ shows that every tail has closure $K$.
\end{proof}

\begin{remark}
The final paragraph is not cosmetic.  Two-sided density alone does not state that arbitrarily late
values see the whole hull.  Compactness of the positive orbit closure supplies the missing
semigroup-to-group step.
\end{remark}

\section{Cofinal values, bounded-gap recurrence, and Haar law}\label{sec:recurrence}

\begin{theorem}[No-gap cofinal value law]\label{thm:no-gap}
Let
\begin{equation}\label{eq:M-def}
 M:=\max_{z\in K}F(z).
\end{equation}
Then, for every $T_0\in\R$,
\begin{equation}\label{eq:tail-closure}
 \cl\{P(t):t\geq T_0\}=F(K)=[0,M].
\end{equation}
For every $y\in[0,M]$ and $\varepsilon>0$, the return set
\begin{equation}\label{eq:approx-return}
 \{t\geq0:|P(t)-y|<\varepsilon\}
\end{equation}
is relatively dense in $[0,\infty)$.  If $0<y<M$, then the exact level set
\begin{equation}\label{eq:exact-level}
 \{t\geq0:P(t)=y\}
\end{equation}
is relatively dense as well.
\end{theorem}

\begin{proof}
By \cref{thm:relation-hull} every tail of $\Theta$ is dense in $K$, and by
\eqref{eq:P-factor} its image has closure $F(K)$.  The latter is compact and connected, hence a
closed interval.  Each summand of $F$ is nonnegative and $F(1)=0$, so the left endpoint is zero.
Compactness gives the maximum $M$, proving \eqref{eq:tail-closure}.

We use a standard compactness form of minimal recurrence.  Let $U\subset K$ be nonempty and open.
For each $x\in K$, density of the positive orbit gives an $s_x\geq0$ such that
$x\Theta(s_x)\in U$.  Therefore the family
\[
 \{U\Theta(-s):s\geq0\}
\]
covers $K$.  Choose a finite subcover with times $s_1,\dots,s_N$, and put
$L=\max_j s_j$.  For any $\tau\geq0$, one has
$\Theta(\tau)\in U\Theta(-s_j)$ for some $j$, and hence
$\Theta(\tau+s_j)\in U$.  Every interval $[\tau,\tau+L]$ therefore contains a return to $U$.

Taking $U=F^{-1}((y-\varepsilon,y+\varepsilon))$ proves relative density of
\eqref{eq:approx-return}.  For $0<y<M$, choose nonempty open sets $U_-$ and $U_+$ on which,
respectively, $F<y$ and $F>y$.  Starting at an arbitrary $\tau$, a return to $U_-$ and a subsequent
return to $U_+$ occur within a uniformly bounded total time.  Continuity of $P$ produces a point
between them with $P(t)=y$.  This proves the claim for \eqref{eq:exact-level}.
\end{proof}

\begin{theorem}[Haar value-distribution without independence]\label{thm:Haar-law}
Let $m_K$ be normalized Haar measure on $K$ and set
\begin{equation}\label{eq:nu-def}
 \nu:=F_*m_K.
\end{equation}
For every bounded continuous $\varphi:\R\to\mathbb C$,
\begin{equation}\label{eq:value-distribution}
 \lim_{T\to\infty}\frac1T\int_0^T\varphi(P(t))\,dt
 =\int_\R\varphi(x)\,d\nu(x).
\end{equation}
Moreover,
\begin{equation}\label{eq:support}
 \supp\nu=[0,M].
\end{equation}
\end{theorem}

\begin{proof}
Let $\chi$ be a character of $K$.  Its restriction to the dense one-parameter subgroup has the
form
\[
 \chi(\Theta(t))=\e^{i\lambda_\chi t}.
\]
If $\chi$ is nontrivial, then $\lambda_\chi\ne0$, since otherwise $\chi$ would be one on a dense
subgroup.  Hence
\[
 \lim_{T\to\infty}\frac1T\int_0^T\chi(\Theta(t))\,dt=0
 =\int_K\chi\,dm_K.
\]
By Stone--Weierstrass, character polynomials are uniformly dense in $C(K)$.  Uniform approximation
therefore gives Haar equidistribution of the flow for every continuous observable.  Apply this to
$\varphi\circ F$ to obtain \eqref{eq:value-distribution}.

The support of $\nu$ is contained in $F(K)$.  If an open interval $V$ meets $F(K)$, then
$F^{-1}(V)$ is a nonempty open subset of $K$.  Every nonempty open subset of a compact group has
positive Haar measure: finitely many translates cover $K$, which would be impossible if its
measure were zero.  Thus $\nu(V)>0$.  Combining this with \cref{thm:no-gap} proves
\eqref{eq:support}.
\end{proof}

The conclusion is deliberately stated for a probability measure rather than a density.  Rational
independence makes $K$ the full product torus and enables product formulas and regularity arguments.
Without independence, $\nu$ may reflect lower-dimensional or infinite-rank relation geometry.  Its
existence and exact support, however, require no independence.

\section{The parity obstruction to the absolute endpoint}

The elementary bound $1-\Re z_i\leq2$ gives $M\leq2A$.  The issue is whether the phases can
simultaneously approach $-1$.  The complete answer depends only on the parity of the relation
lattice.

\begin{theorem}[Universal parity law]\label{thm:parity}
Under \eqref{eq:summable}, the following are equivalent:
\begin{enumerate}
 \item $M=2A$;
 \item the point $(-1)_{i\in I}$ belongs to $K$;
 \item every $k\in\Rel_\omega$ has even total coefficient:
 \begin{equation}\label{eq:even-relation}
    \sum_{i\in I}k_i\equiv0\pmod 2.
 \end{equation}
\end{enumerate}
\end{theorem}

\begin{proof}
For $z\in K$,
\begin{equation}\label{eq:ceiling-defect}
 2A-F(z)=\sum_{i\in I}a_i(1+\Re z_i)\geq0.
\end{equation}
Since all $a_i$ are positive, equality occurs if and only if $\Re z_i=-1$ for every $i$, or
equivalently $z_i=-1$ for every $i$.  The maximum is attained on compact $K$, proving the
equivalence of the first two statements.

By \eqref{eq:annihilator-description}, the all-minus-one point lies in $K$ exactly when
\[
 \prod_i(-1)^{k_i}=(-1)^{\sum_i k_i}=1
\]
for every $k\in\Rel_\omega$.  This is precisely \eqref{eq:even-relation}.
\end{proof}

Rational independence implies $\Rel_\omega=\{0\}$ and hence the endpoint $2A$, but it is much
stronger than necessary.  A relation
\[
 \omega_1+\omega_2=\omega_3+\omega_4
\]
has even total coefficient and creates no parity obstruction.  In contrast,
$\omega_3=\omega_1+\omega_2$ has odd total coefficient and forces $M<2A$.  Thus
\cref{thm:parity} separates relations that genuinely obstruct simultaneous antiphase alignment
from those that do not.

\section{The intrinsic Hilbert screw hull}\label{sec:screw-hull}

The preceding compact group is not an auxiliary device; it is the completion of the associated
screw line in its natural Hilbert metric.  Let $\ell^2(I)$ be the usual complex Hilbert space and
define
\begin{equation}\label{eq:S-map}
 \mathcal S(z)_i:=\sqrt{\frac{a_i}{2}}(1-z_i),
 \qquad z\in K.
\end{equation}
Summability makes this well defined.

\begin{proposition}[Hilbert screw-hull law]\label{prop:Hilbert-hull}
The map $\mathcal S:K\to\ell^2(I)$ is continuous and injective.  For $z,w\in K$,
\begin{equation}\label{eq:hull-metric}
 \|\mathcal S(z)-\mathcal S(w)\|^2
 =F(z\overline w),
\end{equation}
where multiplication and conjugation are coordinatewise.  In particular,
\begin{equation}\label{eq:screw-distance}
 \|\mathcal S(\Theta(t))-\mathcal S(\Theta(u))\|^2=P(t-u).
\end{equation}
The closure of the screw line $t\mapsto\mathcal S(\Theta(t))$ is $\mathcal S(K)$, and the group law
of $K$ transports to a compact connected monothetic group law on this Hilbert-space hull.
\end{proposition}

\begin{proof}
Coordinatewise,
\[
 \frac{a_i}{2}|z_i-w_i|^2
 =a_i\bigl(1-\Re(z_i\overline{w_i})\bigr).
\]
Summation proves \eqref{eq:hull-metric}; \eqref{eq:screw-distance} follows from
\eqref{eq:P-factor}.  Uniform tail control by $2a_i$ proves continuity in the product topology, and
positive weights prove injectivity.  Density of $\Theta(\R)$ gives the closure assertion.
\end{proof}

This gives a geometric reading of the universal results.  The interval $[0,M]$ is the complete
cofinal squared-distance spectrum from the base point.  The parity condition says exactly when the
screw hull contains the coordinatewise antipode.  The Haar law is the canonical value law of the
minimal translation flow on this compact screw hull.

\section{Application to the Riemann screw function}\label{sec:xi}

Assume RH.  Let $\Gamma_+$ be the set of distinct positive ordinates of xi zeros, and let
$m_\gamma$ denote the multiplicity of $\gamma\in\Gamma_+$.  Pairing the positive and negative
zeros in \eqref{eq:suzuki-series} gives
\begin{equation}\label{eq:Psi-positive}
 \Psi(t)=\sum_{\gamma\in\Gamma_+}
  \frac{2m_\gamma}{\gamma^2}\bigl(1-\cos(\gamma t)\bigr).
\end{equation}
The centered Hadamard product gives
\begin{equation}\label{eq:zero-summability}
 \sum_{\gamma\in\Gamma_+}\frac{2m_\gamma}{\gamma^2}<\infty.
\end{equation}
Therefore the universal theorem applies with
$a_\gamma=2m_\gamma/\gamma^2$ and $\omega_\gamma=\gamma$.

Let
\begin{equation}\label{eq:Xi-relations}
 \Rel_\Xi:=\left\{k\in\Z^{(\Gamma_+)}:
        \sum_{\gamma\in\Gamma_+}k_\gamma\gamma=0\right\},
\end{equation}
let $K_\Xi$ be the subgroup of $\T^{\Gamma_+}$ annihilating these relations, and set
\begin{equation}\label{eq:M-Xi}
 M_\Xi:=\max_{z\in K_\Xi}
 \sum_{\gamma\in\Gamma_+}\frac{2m_\gamma}{\gamma^2}(1-\Re z_\gamma).
\end{equation}

\begin{corollary}[Suzuki no-gap recurrence under RH]\label{cor:Suzuki-recurrence}
Assume RH.  No simplicity or rational-independence hypothesis is needed for
\begin{equation}\label{eq:Xi-endpoints}
 \liminf_{t\to\infty}\Psi(t)=0,
 \qquad
 \limsup_{t\to\infty}\Psi(t)=M_\Xi.
\end{equation}
For every $T_0\in\R$,
\begin{equation}\label{eq:Xi-tail-values}
 \cl\{\Psi(t):t\geq T_0\}=[0,M_\Xi].
\end{equation}
Every approximate value in this interval recurs relatively densely, and every interior level is
attained relatively densely.
\end{corollary}

In particular, \eqref{eq:Xi-endpoints} settles the limit-inferior conclusion in
\eqref{eq:suzuki-LI-endpoints} without either of Suzuki's extra assumptions.  The stronger
\eqref{eq:Xi-tail-values} shows that this is not a sparse subsequence phenomenon.

\begin{corollary}[No-LI limiting value law]\label{cor:Suzuki-Haar}
Assume RH.  There is a probability measure $\nu_\Xi$ such that
\begin{equation}\label{eq:Xi-value-law}
 \lim_{T\to\infty}\frac1T\int_0^T\varphi(\Psi(t))\,dt
 =\int_\R\varphi(x)\,d\nu_\Xi(x)
\end{equation}
for every bounded continuous $\varphi$, and
\begin{equation}\label{eq:Xi-measure-support}
 \supp\nu_\Xi=[0,M_\Xi].
\end{equation}
Explicitly, $\nu_\Xi$ is the pushforward of Haar measure on $K_\Xi$ under the observable in
\eqref{eq:M-Xi}.
\end{corollary}

This conclusion complements the $M$-function framework.  Under linear independence the relation
hull is the full product torus, leading to product formulas and, in suitable settings, a continuous
density.  \cref{cor:Suzuki-Haar} keeps existence and determines support for the actual, possibly
relation-constrained zero spectrum.  It does not assert density or absolute continuity.

\begin{corollary}[Parity classification of Suzuki's upper endpoint]\label{cor:Xi-parity}
Assume RH.  Then
\begin{equation}\label{eq:Xi-absolute-ceiling}
 M_\Xi=2\sum_{\gamma}\frac1{\gamma^2}
\end{equation}
if and only if
\begin{equation}\label{eq:Xi-even-relations}
 \sum_{\gamma\in\Gamma_+}k_\gamma\equiv0\pmod2
 \qquad\text{for every }k\in\Rel_\Xi.
\end{equation}
\end{corollary}

\begin{proof}
In the notation of \cref{thm:parity},
\[
 A=\sum_{\gamma\in\Gamma_+}\frac{2m_\gamma}{\gamma^2}
  =\sum_\gamma\frac1{\gamma^2},
\]
where the last sum is over all positive and negative zeros with multiplicity.  Apply
\cref{thm:parity}.
\end{proof}

The condition \eqref{eq:Xi-even-relations} is strictly weaker than linear independence.  It is also
the exact remaining arithmetic obstruction: no generic independence assumption is needed, and no
even-total relation changes Suzuki's absolute endpoint.  We do not know how to prove
\eqref{eq:Xi-even-relations} for the xi zero ordinates.

\begin{corollary}[Almost-periodic RH criterion]\label{cor:AP-RH}
Suzuki's function satisfies
\begin{equation}\label{eq:AP-RH}
 \mathrm{RH}\quad\Longleftrightarrow\quad
 \Psi\text{ is uniformly Bohr almost periodic on }\R.
\end{equation}
\end{corollary}

\begin{proof}
Under RH, \eqref{eq:Psi-positive} converges uniformly and is a uniform limit of trigonometric
polynomials, hence is uniformly Bohr almost periodic.  Conversely, every uniformly almost periodic
function is bounded.  Suzuki's criterion \eqref{eq:suzuki-boundedness} then implies RH.
\end{proof}

\section{Scope and remaining problems}

The compact-hull law answers one problem and isolates another.

It answers the recurrence problem: under RH, zero is a cofinal value of $\Psi$ without simplicity or
LI, indeed with bounded-gap $\varepsilon$-returns; every tail has the same full interval of limit
values; every interior value is crossed with bounded gaps; and the limiting value distribution
exists with exactly that interval as support.  These statements depend only on the summable screw
spectrum and the reality of its frequencies.

It isolates the upper-endpoint problem: Suzuki's displayed absolute endpoint is correct precisely
when the positive xi ordinates have no odd-total integer relation.  Proving this parity assertion is
an arithmetic problem about the zero spectrum and is not supplied by compact dynamics.  It is
weaker than the usual LI conjecture and may therefore be a more accurately targeted conjecture for
the $\limsup$ question.

The results do not prove RH.  They also do not prove regularity of the no-LI limiting measure beyond
its compact support, nor do they establish the cofinal finite-interval operator convergence proposed
in Suzuki's 2026 spectral program.  A natural next question is whether the special positive weights
$2m_\gamma/\gamma^2$ force non-atomicity or absolute continuity of the Haar pushforward under a
relation-rank hypothesis weaker than LI.  A second is whether the compact screw hull can be recovered
as a limit object of Suzuki's finite-aperture self-adjoint realizations.

\begin{thebibliography}{99}

\bibitem{Bohr}
H.~Bohr,
\emph{Almost Periodic Functions},
Chelsea Publishing, New York, 1947.

\bibitem{HewittRoss}
E.~Hewitt and K.~A.~Ross,
\emph{Abstract Harmonic Analysis, Vol.~I},
Springer, Berlin, 1963.

\bibitem{MatsumotoSuzuki2025}
K.~Matsumoto and M.~Suzuki,
\emph{$M$-functions and screw functions: applications to Goldbach's problem and zeros of the
Riemann zeta-function},
arXiv:2409.00888v2, 2025.

\bibitem{Rudin}
W.~Rudin,
\emph{Fourier Analysis on Groups},
Interscience, New York, 1962.

\bibitem{Suzuki2023}
M.~Suzuki,
\emph{Aspects of the screw function corresponding to the Riemann zeta-function},
J. London Math. Soc. \textbf{108} (2023), 1448--1487;
arXiv:2206.03682.

\bibitem{Suzuki2026}
M.~Suzuki,
\emph{Weil's quadratic form via the screw function},
arXiv:2606.09096, 2026.

\end{thebibliography}

\end{document}
